% Options for packages loaded elsewhere
% Options for packages loaded elsewhere
\PassOptionsToPackage{unicode}{hyperref}
\PassOptionsToPackage{hyphens}{url}
%
\documentclass[
  11pt,
  ignorenonframetext,
  aspectratio=169,
]{beamer}
\newif\ifbibliography
\usepackage{pgfpages}
\setbeamertemplate{caption}[numbered]
\setbeamertemplate{caption label separator}{: }
\setbeamercolor{caption name}{fg=normal text.fg}
\beamertemplatenavigationsymbolsempty
% remove section numbering
\setbeamertemplate{part page}{
  \centering
  \begin{beamercolorbox}[sep=16pt,center]{part title}
    \usebeamerfont{part title}\insertpart\par
  \end{beamercolorbox}
}
\setbeamertemplate{section page}{
  \centering
  \begin{beamercolorbox}[sep=12pt,center]{section title}
    \usebeamerfont{section title}\insertsection\par
  \end{beamercolorbox}
}
\setbeamertemplate{subsection page}{
  \centering
  \begin{beamercolorbox}[sep=8pt,center]{subsection title}
    \usebeamerfont{subsection title}\insertsubsection\par
  \end{beamercolorbox}
}
% Prevent slide breaks in the middle of a paragraph
\widowpenalties 1 10000
\raggedbottom
\AtBeginPart{
  \frame{\partpage}
}
\AtBeginSection{
  \ifbibliography
  \else
    \frame{\sectionpage}
  \fi
}
\AtBeginSubsection{
  \frame{\subsectionpage}
}
\usepackage{iftex}
\ifPDFTeX
  \usepackage[T1]{fontenc}
  \usepackage[utf8]{inputenc}
  \usepackage{textcomp} % provide euro and other symbols
\else % if luatex or xetex
  \usepackage{unicode-math} % this also loads fontspec
  \defaultfontfeatures{Scale=MatchLowercase}
  \defaultfontfeatures[\rmfamily]{Ligatures=TeX,Scale=1}
\fi
\usepackage{lmodern}

\ifPDFTeX\else
  % xetex/luatex font selection
\fi
% Use upquote if available, for straight quotes in verbatim environments
\IfFileExists{upquote.sty}{\usepackage{upquote}}{}
\IfFileExists{microtype.sty}{% use microtype if available
  \usepackage[]{microtype}
  \UseMicrotypeSet[protrusion]{basicmath} % disable protrusion for tt fonts
}{}
\makeatletter
\@ifundefined{KOMAClassName}{% if non-KOMA class
  \IfFileExists{parskip.sty}{%
    \usepackage{parskip}
  }{% else
    \setlength{\parindent}{0pt}
    \setlength{\parskip}{6pt plus 2pt minus 1pt}}
}{% if KOMA class
  \KOMAoptions{parskip=half}}
\makeatother

\usepackage{color}
\usepackage{fancyvrb}
\newcommand{\VerbBar}{|}
\newcommand{\VERB}{\Verb[commandchars=\\\{\}]}
\DefineVerbatimEnvironment{Highlighting}{Verbatim}{commandchars=\\\{\}}
% Add ',fontsize=\small' for more characters per line
\usepackage{framed}
\definecolor{shadecolor}{RGB}{241,243,245}
\newenvironment{Shaded}{\begin{snugshade}}{\end{snugshade}}
\newcommand{\AlertTok}[1]{\textcolor[rgb]{0.68,0.00,0.00}{#1}}
\newcommand{\AnnotationTok}[1]{\textcolor[rgb]{0.37,0.37,0.37}{#1}}
\newcommand{\AttributeTok}[1]{\textcolor[rgb]{0.40,0.45,0.13}{#1}}
\newcommand{\BaseNTok}[1]{\textcolor[rgb]{0.68,0.00,0.00}{#1}}
\newcommand{\BuiltInTok}[1]{\textcolor[rgb]{0.00,0.23,0.31}{#1}}
\newcommand{\CharTok}[1]{\textcolor[rgb]{0.13,0.47,0.30}{#1}}
\newcommand{\CommentTok}[1]{\textcolor[rgb]{0.37,0.37,0.37}{#1}}
\newcommand{\CommentVarTok}[1]{\textcolor[rgb]{0.37,0.37,0.37}{\textit{#1}}}
\newcommand{\ConstantTok}[1]{\textcolor[rgb]{0.56,0.35,0.01}{#1}}
\newcommand{\ControlFlowTok}[1]{\textcolor[rgb]{0.00,0.23,0.31}{\textbf{#1}}}
\newcommand{\DataTypeTok}[1]{\textcolor[rgb]{0.68,0.00,0.00}{#1}}
\newcommand{\DecValTok}[1]{\textcolor[rgb]{0.68,0.00,0.00}{#1}}
\newcommand{\DocumentationTok}[1]{\textcolor[rgb]{0.37,0.37,0.37}{\textit{#1}}}
\newcommand{\ErrorTok}[1]{\textcolor[rgb]{0.68,0.00,0.00}{#1}}
\newcommand{\ExtensionTok}[1]{\textcolor[rgb]{0.00,0.23,0.31}{#1}}
\newcommand{\FloatTok}[1]{\textcolor[rgb]{0.68,0.00,0.00}{#1}}
\newcommand{\FunctionTok}[1]{\textcolor[rgb]{0.28,0.35,0.67}{#1}}
\newcommand{\ImportTok}[1]{\textcolor[rgb]{0.00,0.46,0.62}{#1}}
\newcommand{\InformationTok}[1]{\textcolor[rgb]{0.37,0.37,0.37}{#1}}
\newcommand{\KeywordTok}[1]{\textcolor[rgb]{0.00,0.23,0.31}{\textbf{#1}}}
\newcommand{\NormalTok}[1]{\textcolor[rgb]{0.00,0.23,0.31}{#1}}
\newcommand{\OperatorTok}[1]{\textcolor[rgb]{0.37,0.37,0.37}{#1}}
\newcommand{\OtherTok}[1]{\textcolor[rgb]{0.00,0.23,0.31}{#1}}
\newcommand{\PreprocessorTok}[1]{\textcolor[rgb]{0.68,0.00,0.00}{#1}}
\newcommand{\RegionMarkerTok}[1]{\textcolor[rgb]{0.00,0.23,0.31}{#1}}
\newcommand{\SpecialCharTok}[1]{\textcolor[rgb]{0.37,0.37,0.37}{#1}}
\newcommand{\SpecialStringTok}[1]{\textcolor[rgb]{0.13,0.47,0.30}{#1}}
\newcommand{\StringTok}[1]{\textcolor[rgb]{0.13,0.47,0.30}{#1}}
\newcommand{\VariableTok}[1]{\textcolor[rgb]{0.07,0.07,0.07}{#1}}
\newcommand{\VerbatimStringTok}[1]{\textcolor[rgb]{0.13,0.47,0.30}{#1}}
\newcommand{\WarningTok}[1]{\textcolor[rgb]{0.37,0.37,0.37}{\textit{#1}}}

\usepackage{longtable,booktabs,array}
\newcounter{none} % for unnumbered tables
\usepackage{calc} % for calculating minipage widths
\usepackage{caption}
% Make caption package work with longtable
\makeatletter
\def\fnum@table{\tablename~\thetable}
\makeatother
\usepackage{graphicx}
\makeatletter
\newsavebox\pandoc@box
\newcommand*\pandocbounded[1]{% scales image to fit in text height/width
  \sbox\pandoc@box{#1}%
  \Gscale@div\@tempa{\textheight}{\dimexpr\ht\pandoc@box+\dp\pandoc@box\relax}%
  \Gscale@div\@tempb{\linewidth}{\wd\pandoc@box}%
  \ifdim\@tempb\p@<\@tempa\p@\let\@tempa\@tempb\fi% select the smaller of both
  \ifdim\@tempa\p@<\p@\scalebox{\@tempa}{\usebox\pandoc@box}%
  \else\usebox{\pandoc@box}%
  \fi%
}
% Set default figure placement to htbp
\def\fps@figure{htbp}
\makeatother





\setlength{\emergencystretch}{3em} % prevent overfull lines

\providecommand{\tightlist}{%
  \setlength{\itemsep}{0pt}\setlength{\parskip}{0pt}}



 


\setbeamertemplate{navigation symbols}{}
\setbeamertemplate{footline}[frame number]
\setbeamertemplate{itemize items}[circle]
\setbeamertemplate{itemize subitem}[triangle]
\setbeamercolor{frametitle}{fg=black}
\setbeamerfont{frametitle}{series=\bfseries}
\usepackage{amsmath}
\usepackage{booktabs}
\usepackage{array}
\newcolumntype{R}[1]{>{\raggedright\arraybackslash}p{#1}}
\usepackage{fancyvrb}
\fvset{fontsize=\scriptsize}
\makeatletter
\def\verbatim@font{\scriptsize\ttfamily}
\makeatother
\makeatletter
\@ifpackageloaded{caption}{}{\usepackage{caption}}
\AtBeginDocument{%
\ifdefined\contentsname
  \renewcommand*\contentsname{Table of contents}
\else
  \newcommand\contentsname{Table of contents}
\fi
\ifdefined\listfigurename
  \renewcommand*\listfigurename{List of Figures}
\else
  \newcommand\listfigurename{List of Figures}
\fi
\ifdefined\listtablename
  \renewcommand*\listtablename{List of Tables}
\else
  \newcommand\listtablename{List of Tables}
\fi
\ifdefined\figurename
  \renewcommand*\figurename{Figure}
\else
  \newcommand\figurename{Figure}
\fi
\ifdefined\tablename
  \renewcommand*\tablename{Table}
\else
  \newcommand\tablename{Table}
\fi
}
\@ifpackageloaded{float}{}{\usepackage{float}}
\floatstyle{ruled}
\@ifundefined{c@chapter}{\newfloat{codelisting}{h}{lop}}{\newfloat{codelisting}{h}{lop}[chapter]}
\floatname{codelisting}{Listing}
\newcommand*\listoflistings{\listof{codelisting}{List of Listings}}
\makeatother
\makeatletter
\makeatother
\makeatletter
\@ifpackageloaded{caption}{}{\usepackage{caption}}
\@ifpackageloaded{subcaption}{}{\usepackage{subcaption}}
\makeatother

\usepackage{bookmark}
\IfFileExists{xurl.sty}{\usepackage{xurl}}{} % add URL line breaks if available
\urlstyle{same}
\hypersetup{
  pdftitle={Difference-in-Differences: Fundamentals},
  pdfauthor={Professor Ji-Woong Chung},
  hidelinks,
  pdfcreator={LaTeX via pandoc}}


\title{\texorpdfstring{Difference-in-Differences:
Fundamentals}{Difference-in-Differences: Fundamentals}}
\subtitle{\texorpdfstring{BUSS975: Causal Inference in Financial
Research}{BUSS975: Causal Inference in Financial Research}}
\author{Professor Ji-Woong Chung}
\date{}
\institute{Korea University}

\begin{document}
\frame{\titlepage}


\begin{frame}{Where we left off}
\protect\phantomsection\label{where-we-left-off}
\begin{itemize}
\tightlist
\item
  Lecture 6a ended on a \(2 \times 2\) panel: two firms, two years, one
  of them treated in year two. Firm effects removed \emph{size}, year
  effects removed the \emph{boom}, and what was left was
  \((15 - 10) - (7 - 4) = 2\).
\item
  That calculation has a name. It is older than fixed effects, older
  than regression, and was invented three separate times by people
  trying to save lives.
\end{itemize}

\vspace{0.4em}

\begin{itemize}
\tightlist
\item
  Lecture 6b asked of every estimator: \emph{which way is it wrong?} The
  answer there came from the model.
\item
  This week the answer comes from one \textbf{counterfactual} you will
  never see: what would have happened to the treated, had they not been
  treated. Everything in difference-in-differences is a claim about that
  one number.
\end{itemize}
\end{frame}

\begin{frame}{The plan}
\protect\phantomsection\label{the-plan}
\begin{enumerate}
\tightlist
\item
  \textbf{Three births.} Vienna, London, Washington --- and why the
  method keeps being reinvented.
\item
  \textbf{Four averages and three subtractions.} The \(2 \times 2\), and
  three regressions that compute it exactly.
\item
  \textbf{What makes it causal.} Potential outcomes, parallel trends, no
  anticipation, and what goes wrong without them.
\item
  \textbf{Card and Krueger.} The study that made it mainstream, and a
  picture of what regression imputes.
\item
  \textbf{Event studies and falsifications.} Evidence for an assumption
  you cannot test.
\item
  \textbf{Who gets treated, and why.} Which assignment mechanisms break
  parallel trends --- and which, surprisingly, do not.
\end{enumerate}

\vspace{0.4em}

\begin{itemize}
\tightlist
\item
  One running question: \textbf{what is the counterfactual trend, and
  why do I believe the control group has it?}
\end{itemize}
\end{frame}

\section{1. Three Births}\label{three-births}

\begin{frame}{Vienna, 1847: Semmelweis}
\protect\phantomsection\label{vienna-1847-semmelweis}
\begin{center}
\includegraphics[width=0.8\linewidth,height=\textheight,keepaspectratio]{07a_did_fundamentals_files/figure-beamer/semmelweis-1.pdf}
\end{center}

\begin{itemize}
\tightlist
\item
  Two maternity clinics in the same hospital. Women were assigned by
  \textbf{arrival time} --- by law, so effectively at random. Yet
  mortality in the First Clinic ran three to four times the Second's.
\item
  Semmelweis's guess: doctors came to the delivery room straight from
  the autopsy room. In May 1847 he made them wash their hands in
  chlorine. \textbf{The First Clinic fell to the Second Clinic's level
  within a year.}
\item
  Two groups, before and after, one of them treated. Nobody believed
  him. He died in an asylum.
\end{itemize}
\end{frame}

\begin{frame}{London, 1854: Snow's grand experiment}
\protect\phantomsection\label{london-1854-snows-grand-experiment}
\begin{itemize}
\tightlist
\item
  Cholera killed tens of thousands in London. The medical consensus was
  \emph{miasma}: disease travels in foul air. Snow thought it travelled
  in \textbf{water}.
\item
  Two companies supplied the same South London streets --- sometimes the
  same street, house by house. Between 1849 and 1854 \textbf{Lambeth}
  moved its intake upstream of the sewage outfalls; \textbf{Southwark
  and Vauxhall} did not.
\item
  Snow went door to door asking which company served each house. Where
  nobody knew, he took a water sample: Lambeth's upstream water had less
  salt. \emph{Shoe leather and design, together.}
\end{itemize}

\vspace{0.3em}

\begin{itemize}
\tightlist
\item
  What he had: cholera mortality for each company's customers, in 1849
  \textbf{and} 1854. Two groups, two periods, one of them changed in
  between.
\item
  He was not believed either. Miasma was too entrenched; the data could
  not crack it.
\end{itemize}
\end{frame}

\begin{frame}{Snow, as four averages and three subtractions}
\protect\phantomsection\label{snow-as-four-averages-and-three-subtractions}
Write mortality as a sum of pieces. \(L\) and \(SV\) are the companies'
levels in 1849; \(L_t\) and \(SV_t\) are how each would have moved by
1854 \emph{on its own}; \(\delta\) is the effect of the pipe.

\begin{center}\small
\begin{tabular}{llccc}
\toprule
Company & Year & Mortality & First difference & Second difference \\
\midrule
Lambeth & 1849 & $L$ & & \\
        & 1854 & $L + L_t + \delta$ & $L_t + \delta$ & \\
        &      &                    &                & $\delta + (L_t - SV_t)$ \\
Southwark \& Vauxhall & 1849 & $SV$ & & \\
        & 1854 & $SV + SV_t$ & $SV_t$ & \\
\bottomrule
\end{tabular}
\end{center}

\begin{itemize}
\tightlist
\item
  \textbf{Subtract within each group}: levels cancel. Lambeth is left
  with its trend plus the effect; S\&V with its trend alone.
\item
  \textbf{Subtract across groups}: trends cancel --- \emph{if}
  \(L_t = SV_t\). That equality is the whole method.
\item
  Snow's numbers (per 10,000 houses): Lambeth \(62 \to 14\), S\&V
  \(565 \to 614\). So \(\hat\delta = (-48) - (+49) = \mathbf{-97}\),
  provided Lambeth's mortality would otherwise have risen like S\&V's.
\end{itemize}
\end{frame}

\begin{frame}{Washington, 1970s: Ashenfelter}
\protect\phantomsection\label{washington-1970s-ashenfelter}
\begin{itemize}
\tightlist
\item
  Orley Ashenfelter, a Princeton labour economist, went to the
  Department of Labor to evaluate \textbf{job training programmes}. He
  had a rare thing: longitudinal earnings data on trainees and
  non-trainees.
\item
  He estimated what we now call two-way fixed effects --- thousands of
  individual effects plus time effects --- and had to explain the result
  to \textbf{people who did not want to hear the word ``regression.''}
\item
  So he described the coefficient the way it actually computes: four
  averages and three subtractions. He called it
  \textbf{difference-in-differences.}
\end{itemize}

\vspace{0.3em}

\begin{itemize}
\tightlist
\item
  He was not dumbing it down. We will see in Section 2 that the
  regression coefficient \emph{is} that arithmetic, to the last digit.
\item
  Three births, one trait: a method that could be \textbf{shown to
  someone with no training}, and whose weak point --- the counterfactual
  trend --- is visible on the page.
\end{itemize}
\end{frame}

\section{2. Four Averages and Three
Subtractions}\label{four-averages-and-three-subtractions}

\begin{frame}{The \(2 \times 2\) in one line}
\protect\phantomsection\label{the-2-times-2-in-one-line}
Let \(k\) be the group treated at one date, \(U\) a group never treated
in the window, and \emph{pre} / \emph{post} the periods on either side:
\[
\hat\delta_{DiD} \;=\; \big(\bar Y_k^{\,post} - \bar Y_k^{\,pre}\big) \;-\; \big(\bar Y_U^{\,post} - \bar Y_U^{\,pre}\big)
\]

\begin{itemize}
\tightlist
\item
  \textbf{Four averages}: each group in each period. \textbf{Three
  subtractions}: two within groups, one across.
\item
  The literature now calls this \textbf{the \(2 \times 2\)}
  (Goodman-Bacon 2021; Baker et al.~2025). Everything more complex is
  built from it, so it is worth knowing in your sleep.
\item
  The within-group subtractions delete anything \textbf{fixed about the
  group}. The across-group subtraction deletes anything \textbf{common
  to the period}. What survives is the effect --- and any difference in
  trends.
\end{itemize}
\end{frame}

\begin{frame}{The same number from a regression}
\protect\phantomsection\label{the-same-number-from-a-regression}
\[
Y_{it} \;=\; \alpha_0 + \alpha_1\,\text{Treat}_i + \alpha_2\,\text{Post}_t + \delta\,(\text{Treat}_i \times \text{Post}_t) + \varepsilon_{it}
\]

Four dummy cells, four parameters: OLS fits each cell mean exactly.

\begin{center}\small
\begin{tabular}{lcc}
\toprule
 & Pre & Post \\
\midrule
Control ($\text{Treat}=0$) & $\hat\alpha_0$ & $\hat\alpha_0 + \hat\alpha_2$ \\
Treated ($\text{Treat}=1$) & $\hat\alpha_0 + \hat\alpha_1$ & $\hat\alpha_0 + \hat\alpha_1 + \hat\alpha_2 + \hat\delta$ \\
\bottomrule
\end{tabular}
\end{center}

Plug the cells into the \(2 \times 2\): \[
\hat\delta_{DiD} = \big[(\hat\alpha_0 + \hat\alpha_1 + \hat\alpha_2 + \hat\delta) - (\hat\alpha_0 + \hat\alpha_1)\big] - \big[(\hat\alpha_0 + \hat\alpha_2) - \hat\alpha_0\big] = (\hat\alpha_2 + \hat\delta) - \hat\alpha_2 = \hat\delta
\]

\begin{itemize}
\tightlist
\item
  Not approximately. \textbf{Identically.} The interaction coefficient
  is four averages and three subtractions.
\end{itemize}
\end{frame}

\begin{frame}{Castle doctrine: the data}
\protect\phantomsection\label{castle-doctrine-the-data}
Cheng and Hoekstra (2013): a balanced panel of \(50\) states,
\(2000\)--\(2010\). Over the decade, \(21\) states passed
\textbf{``stand your ground''} laws extending the right to use lethal
force in self-defence beyond the home. Outcome: \textbf{log homicides}
per \(100{,}000\).

\begin{itemize}
\tightlist
\item
  To keep it \(2 \times 2\): keep the \(13\) states that adopted in
  \textbf{2006} and the \(29\) that never adopted; drop the other
  cohorts. Pre \(= 2000\)--\(05\), post \(= 2006\)--\(10\).
\end{itemize}

\begin{center}\small
\begin{tabular}{lccc}
\toprule
Mean log homicide & Pre & Post & Change \\
\midrule
Never adopted & $1.252$ & $1.196$ & $-0.055$ \\
Adopted in 2006 & $1.786$ & $1.799$ & $0.013$ \\
\midrule
Difference-in-differences & & & $\mathbf{0.068}$ \\
\bottomrule
\end{tabular}
\end{center}

\begin{itemize}
\tightlist
\item
  Adopters had \textbf{higher} homicide rates all along: the level
  difference is not the effect. Homicide fell in non-adopting states and
  rose in adopting ones; the gap between those changes is the
  \(2 \times 2\): about \(7\) percent more homicide.
\end{itemize}
\end{frame}

\begin{frame}[fragile]{Four ways to the same number}
\protect\phantomsection\label{four-ways-to-the-same-number}
\begin{Shaded}
\begin{Highlighting}[]
\NormalTok{k }\OtherTok{\textless{}{-}} \FunctionTok{subset}\NormalTok{(castle, }\FunctionTok{is.na}\NormalTok{(effyear) }\SpecialCharTok{|}\NormalTok{ effyear }\SpecialCharTok{==} \DecValTok{2006}\NormalTok{)          }\CommentTok{\# 2006 cohort + never treated}
\NormalTok{k}\SpecialCharTok{$}\NormalTok{treat }\OtherTok{\textless{}{-}} \SpecialCharTok{!}\FunctionTok{is.na}\NormalTok{(k}\SpecialCharTok{$}\NormalTok{effyear); k}\SpecialCharTok{$}\NormalTok{post }\OtherTok{\textless{}{-}}\NormalTok{ k}\SpecialCharTok{$}\NormalTok{year }\SpecialCharTok{\textgreater{}=} \DecValTok{2006}
\NormalTok{m }\OtherTok{\textless{}{-}} \FunctionTok{with}\NormalTok{(k, }\FunctionTok{tapply}\NormalTok{(l\_homicide, }\FunctionTok{list}\NormalTok{(treat, post), mean))       }\CommentTok{\# 1. the 2x2}
\NormalTok{(m[}\DecValTok{2}\NormalTok{, }\DecValTok{2}\NormalTok{] }\SpecialCharTok{{-}}\NormalTok{ m[}\DecValTok{2}\NormalTok{, }\DecValTok{1}\NormalTok{]) }\SpecialCharTok{{-}}\NormalTok{ (m[}\DecValTok{1}\NormalTok{, }\DecValTok{2}\NormalTok{] }\SpecialCharTok{{-}}\NormalTok{ m[}\DecValTok{1}\NormalTok{, }\DecValTok{1}\NormalTok{])}
\FunctionTok{feols}\NormalTok{(l\_homicide }\SpecialCharTok{\textasciitilde{}}\NormalTok{ treat }\SpecialCharTok{*}\NormalTok{ post, k, }\AttributeTok{cluster =} \SpecialCharTok{\textasciitilde{}}\NormalTok{sid)            }\CommentTok{\# 2. interaction OLS}
\FunctionTok{feols}\NormalTok{(l\_homicide }\SpecialCharTok{\textasciitilde{}}\NormalTok{ treat}\SpecialCharTok{:}\NormalTok{post }\SpecialCharTok{|}\NormalTok{ sid }\SpecialCharTok{+}\NormalTok{ year, k, }\AttributeTok{cluster =} \SpecialCharTok{\textasciitilde{}}\NormalTok{sid) }\CommentTok{\# 3. two{-}way fixed effects}
\NormalTok{w }\OtherTok{\textless{}{-}}\NormalTok{ ...  }\CommentTok{\# each state\textquotesingle{}s post{-}period mean minus its pre{-}period mean}
\FunctionTok{feols}\NormalTok{(dbar }\SpecialCharTok{\textasciitilde{}}\NormalTok{ treat, w, }\AttributeTok{cluster =} \SpecialCharTok{\textasciitilde{}}\NormalTok{sid)                          }\CommentTok{\# 4. the long difference}
\end{Highlighting}
\end{Shaded}

\begin{center}\small
\begin{tabular}{lcccc}
\toprule
 & $2 \times 2$ & Interaction OLS & TWFE & Long difference \\
\midrule
$\hat\delta$ & $0.0682359$ & $0.0682359$ & $0.0682359$ & $0.0682359$ \\
\bottomrule
\end{tabular}
\end{center}

\begin{itemize}
\tightlist
\item
  \textbf{Three regressions, one arithmetic.} With two groups and two
  periods, unit and year fixed effects are just \(\text{Treat}_i\) and
  \(\text{Post}_t\) in disguise, and the long difference \emph{is} the
  within-group subtraction.
\item
  This is why ``difference-in-differences'' and ``two-way fixed
  effects'' were synonyms for thirty years. They stop being synonyms
  when treatment dates differ across units --- next lecture.
\end{itemize}
\end{frame}

\begin{frame}{Why regression, then?}
\protect\phantomsection\label{why-regression-then}
\begin{itemize}
\tightlist
\item
  Not for the point estimate: that is fixed by the arithmetic.
\item
  For \textbf{standard errors} --- and for adding covariates, more
  periods, more groups. All of which, we will see, come with strings
  attached.
\end{itemize}

\vspace{0.4em}

\begin{itemize}
\tightlist
\item
  Ashenfelter's lesson is about \textbf{communication}. The regression
  and the four averages are the same object; one of them can be
  explained to a minister, a journalist, a referee outside the field.
\item
  Keep the arithmetic in reach. When a sophisticated estimator gives a
  number you cannot reproduce from a few group means, you do not yet
  understand what it did.
\end{itemize}
\end{frame}

\section{3. What Makes It Causal}\label{what-makes-it-causal}

\begin{frame}{From subtraction to causality: potential outcomes}
\protect\phantomsection\label{from-subtraction-to-causality-potential-outcomes}
Write the \(2 \times 2\) with expectations and swap in potential
outcomes using the switching equation \(Y = D\,Y^1 + (1-D)\,Y^0\): \[
\hat\delta_{DiD} = \big(\mathbb{E}[Y_k|\text{Post}] - \mathbb{E}[Y_k|\text{Pre}]\big) - \big(\mathbb{E}[Y_U|\text{Post}] - \mathbb{E}[Y_U|\text{Pre}]\big)
\] \[
= \big(\mathbb{E}[Y^{1}_k|\text{Post}] - \mathbb{E}[Y^{0}_k|\text{Pre}]\big) - \big(\mathbb{E}[Y^{0}_U|\text{Post}] - \mathbb{E}[Y^{0}_U|\text{Pre}]\big)
\]

\begin{itemize}
\tightlist
\item
  Only one of the four cells is treated: \textbf{group \(k\), after.}
  Every other cell is \(Y^0\).
\item
  Nothing here is a causal effect yet. A causal effect is \(Y^1 - Y^0\)
  \textbf{for the same group at the same time}. Each pair above compares
  either two groups, or two dates.
\end{itemize}
\end{frame}

\begin{frame}{Adding a zero}
\protect\phantomsection\label{adding-a-zero}
Add and subtract the one thing you cannot observe --- the treated
group's untreated outcome after treatment,
\(\mathbb{E}[Y^0_k|\text{Post}]\) --- and regroup: \begin{align*}
\hat\delta_{DiD} \;=\;& \underbrace{\mathbb{E}[Y^1_k|\text{Post}] - \mathbb{E}[Y^0_k|\text{Post}]}_{\textbf{ATT}} \\
&+\; \underbrace{\big(\mathbb{E}[Y^0_k|\text{Post}] - \mathbb{E}[Y^0_k|\text{Pre}]\big) - \big(\mathbb{E}[Y^0_U|\text{Post}] - \mathbb{E}[Y^0_U|\text{Pre}]\big)}_{\textbf{non-parallel trends bias}}
\end{align*}

\begin{itemize}
\tightlist
\item
  \textbf{The first term is the only causal parameter DiD can deliver}:
  the average effect \emph{on the treated}, \emph{in the post period}.
  Say it that precisely every time.
\item
  \textbf{The second term is a selection bias} --- not in levels of
  \(Y^0\), as in Lecture 3, but in \emph{changes} of \(Y^0\). It is
  itself a \(2 \times 2\), computed on \(Y^0\) instead of \(Y\).
\item
  Parallel trends is the assumption that the second term is zero: \[
  \mathbb{E}[Y^0_k|\text{Post}] - \mathbb{E}[Y^0_k|\text{Pre}] \;=\; \mathbb{E}[Y^0_U|\text{Post}] - \mathbb{E}[Y^0_U|\text{Pre}]
  \]
\end{itemize}
\end{frame}

\begin{frame}{What parallel trends says, and does not}
\protect\phantomsection\label{what-parallel-trends-says-and-does-not}
\begin{itemize}
\tightlist
\item
  It is a statement about a \textbf{counterfactual}: how the treated
  group's \(Y^0\) \emph{would have} moved. One of its four terms does
  not exist in any dataset. \textbf{It cannot be tested.} It can only be
  argued.
\item
  It says nothing about \textbf{levels}. Lambeth and S\&V could differ
  tenfold in 1849; adopters can have twice the homicide rate. Fixed
  differences are subtracted out. What must match is the \emph{change}.
\item
  It is \textbf{specific to the comparison and the baseline you chose.}
  Change the control group, or the pre period, and you have made a
  different assumption.
\item
  It is a statement in a particular \textbf{scale}: parallel in levels
  is not parallel in logs. Pick the scale by the economics, before you
  see the estimate.
\end{itemize}

\vspace{0.3em}

\begin{itemize}
\tightlist
\item
  Three assumptions, in total: \textbf{parallel trends}, \textbf{no
  anticipation} (next), and SUTVA, which is the same as always --- one
  version of the treatment, no spillovers onto the controls.
\end{itemize}
\end{frame}

\begin{frame}{A numerical example: parallel trends holds}
\protect\phantomsection\label{a-numerical-example-parallel-trends-holds}
Two groups, \(1980\)--\(1990\), group \(1\) treated from \(1986\). The
simulation knows both potential outcomes; the data show only one.

\begin{center}\small
\begin{tabular}{lcc@{\hspace{2.5em}}cc}
\toprule
 & \multicolumn{2}{c}{Observed $Y$} & \multicolumn{2}{c}{Untreated $Y^0$} \\
Average & Group 1 & Group 2 & Group 1 & Group 2 \\
\midrule
Pre (1980--85)  & $6.06$  & $6.07$  & $6.06$  & $6.07$ \\
Post (1986--90) & $41.56$ & $11.59$ & $11.57$ & $11.59$ \\
Change          & $35.51$ & $5.52$  & $5.51$  & $5.52$ \\
\bottomrule
\end{tabular}
\end{center}

\begin{itemize}
\tightlist
\item
  \textbf{DiD} \(= 35.51 - 5.52 = 29.99\), from the observed columns.
\item
  \textbf{ATT} \(= 41.56 - 11.57 = 30.0\): group 1's observed post
  outcome minus its \emph{untreated} post outcome.
\item
  \textbf{Non-parallel trends bias} \(= 5.51 - 5.52 = -0.01\), the
  \(2 \times 2\) on \(Y^0\). And \(29.99 = 30.0 + (-0.01)\).
\end{itemize}
\end{frame}

\begin{frame}{Same data, different counterfactual}
\protect\phantomsection\label{same-data-different-counterfactual}
Change nothing observable. Only group 1's post-period \(Y^0\) --- the
column no one sees --- is different.

\begin{center}\small
\begin{tabular}{lcc@{\hspace{2.5em}}cc}
\toprule
 & \multicolumn{2}{c}{Observed $Y$} & \multicolumn{2}{c}{Untreated $Y^0$} \\
Average & Group 1 & Group 2 & Group 1 & Group 2 \\
\midrule
Pre (1980--85)  & $6.06$  & $6.07$  & $6.06$  & $6.07$ \\
Post (1986--90) & $41.56$ & $11.59$ & $\mathbf{36.60}$ & $11.59$ \\
Change          & $35.51$ & $5.52$  & $\mathbf{30.55}$ & $5.52$ \\
\bottomrule
\end{tabular}
\end{center}

\begin{itemize}
\tightlist
\item
  \textbf{DiD} \(= 29.99\). Unchanged --- it uses no counterfactual.
\item
  \textbf{ATT} \(= 41.56 - 36.60 = 4.96\). \textbf{Bias}
  \(= 30.55 - 5.52 = 25.03\). Group 1 was going to rise anyway.
\item
  \textbf{The point.} Two worlds, identical data, same \(\hat\delta\),
  ATT of \(30\) in one and \(5\) in the other. The data cannot tell you
  which world you are in. Only an argument about \emph{why} group 1 was
  treated can --- Section 6.
\end{itemize}
\end{frame}

\begin{frame}{No anticipation: what it literally says}
\protect\phantomsection\label{no-anticipation-what-it-literally-says}
\begin{itemize}
\item
  The name suggests psychology. The content is mundane: \textbf{the
  treated group is untreated until it is treated} --- in the pre period,
  \(Y = Y^0\) for group \(k\).
\item
  Suppose it fails: the ``pre'' period is already treated, so
  \(Y_k^{pre} = Y^1_k\). Add two zeros and regroup: \[
  \hat\delta_{DiD} \;=\; \underbrace{\text{ATT}_k(\text{Post})}_{\text{what you want}} \;+\; \underbrace{\text{PT bias}}_{\text{as before}} \;-\; \underbrace{\text{ATT}_k(\text{Pre})}_{\text{the effect already in the baseline}}
  \]
\item
  With \textbf{constant} effects,
  \(\text{ATT}(\text{Post}) = \text{ATT}(\text{Pre})\) and DiD is
  \textbf{zero}. With effects that \textbf{grow}, DiD is the
  \emph{growth} in the effect, not the effect.
\item
  Why the name: if units react to a policy \emph{before} it is enforced
  --- firms hire, fire, exit in anticipation --- then the baseline is
  already treated, and this is what you estimate.
\end{itemize}
\end{frame}

\begin{frame}{No anticipation, simulated}
\protect\phantomsection\label{no-anticipation-simulated}
\(1{,}000\) firms, \(1990\)--\(93\), half treated from \(1991\) onward.
\(Y^0\) has a firm effect and a common trend. Two effect paths:
\textbf{constant} (\(10, 10, 10\)) and \textbf{dynamic}
(\(10, 20, 30\)).

\begin{center}\small
\begin{tabular}{lcccc}
\toprule
 & \multicolumn{2}{c}{Baseline $= 1990$ (clean)} & \multicolumn{2}{c}{Baseline $= 1991$ (already treated)} \\
 & constant & dynamic & constant & dynamic \\
\midrule
$\hat\delta_{DiD}$ & $10.03$ & $20.03$ & $0.06$ & $15.06$ \\
ATT over the post period & $10$ & $20$ & $10$ & $25$ \\
\bottomrule
\end{tabular}
\end{center}

\begin{itemize}
\tightlist
\item
  Clean baseline: DiD hits the ATT in both columns.
\item
  Treated baseline, constant effect: \textbf{zero}. The effect is in
  both periods and cancels.
\item
  Treated baseline, dynamic effect: \(25 - 10 = 15\). You estimate the
  \emph{change} in the effect and call it the effect.
\end{itemize}
\end{frame}

\begin{frame}{Dating the treatment}
\protect\phantomsection\label{dating-the-treatment}
\begin{itemize}
\tightlist
\item
  \textbf{The baseline must be clean.} If a group was treated even
  briefly in a period, that period belongs to \emph{post}. Err on the
  side of starting treatment early.
\item
  \textbf{Announcement vs.~enforcement.} If units plausibly respond when
  a law is \emph{passed} rather than when it takes effect, make the
  baseline the period before passage.
\end{itemize}

\vspace{0.3em}

\begin{itemize}
\tightlist
\item
  Be aware of what that choice changes:

  \begin{itemize}
  \tightlist
  \item
    the \textbf{ATT} --- it now averages over an ``announced but not
    enforced'' period as well;
  \item
    the \textbf{parallel trends assumption} --- it is always relative to
    a baseline, and you just moved it.
  \end{itemize}
\item
  None of this is right or wrong in the abstract. It is a design
  decision; make it on purpose, and say what you did.
\end{itemize}
\end{frame}

\begin{frame}{Avoid already-treated controls}
\protect\phantomsection\label{avoid-already-treated-controls}
Not an assumption --- a strong recommendation. Suppose the control group
\(U\) was treated \emph{before} the window, so \(Y_U = Y^1_U\) in both
periods. Add three zeros and regroup: \[
\hat\delta_{DiD} \;=\; \text{ATT}_k \;+\; \text{PT bias} \;-\; \underbrace{\big[\text{ATT}_U(\text{Post}) - \text{ATT}_U(\text{Pre})\big]}_{\Delta \text{ATT}_U:\ \text{how the \emph{control's} effect changed}}
\]

\begin{itemize}
\tightlist
\item
  If the control's effect is \textbf{constant}, no harm: it differences
  out like a fixed effect.
\item
  If it is \textbf{still evolving} --- fading, growing --- that
  evolution is subtracted from your estimate and looks like a treatment
  effect of the opposite sign.
\item
  Obvious? Nobody would \emph{choose} an already-treated control. But
  when treatment dates differ across units, standard two-way fixed
  effects \textbf{makes these comparisons for you, silently.} That is
  the next lecture's subject.
\end{itemize}
\end{frame}

\section{4. Card and Krueger}\label{card-and-krueger}

\begin{frame}{New Jersey vs.~Pennsylvania, 1992}
\protect\phantomsection\label{new-jersey-vs.-pennsylvania-1992}
\begin{itemize}
\tightlist
\item
  Late 1991: New Jersey will raise its minimum wage from \$4.25 to
  \$5.05 in April 1992. Pennsylvania, next door, will not.
\item
  Card and Krueger phoned every fast-food restaurant they could find in
  both states --- \textbf{February} (before) and \textbf{November}
  (after) --- and asked about wages and employment. They ran their own
  survey because no dataset existed.
\end{itemize}

\vspace{0.3em}

\begin{itemize}
\tightlist
\item
  \textbf{First, show the policy had bite.} In November, the NJ wage
  distribution has a spike at exactly \$5.05: a large mass of workers
  was pushed up to the new minimum. In PA, no spike. The treatment
  \emph{happened} to the treated.
\item
  Always do this before looking at the outcome. If the policy did not
  move the thing it was meant to move, there is no reason for anything
  downstream to have moved either.
\end{itemize}
\end{frame}

\begin{frame}{The result: four averages, three subtractions}
\protect\phantomsection\label{the-result-four-averages-three-subtractions}
Full-time-equivalent employment per store (Card and Krueger 1994, Table
3):

\begin{center}\small
\begin{tabular}{lccc}
\toprule
 & PA & NJ & NJ $-$ PA \\
\midrule
FTE, February (before) & $23.33$ \,($1.35$) & $20.44$ \,($0.51$) & $-2.89$ \,($1.44$) \\
FTE, November (after)  & $21.17$ \,($0.94$) & $21.03$ \,($0.52$) & $-0.14$ \,($1.07$) \\
Change                 & $-2.16$ \,($1.25$) & $+0.59$ \,($0.54$) & $\mathbf{+2.76}$ \,($1.36$) \\
\bottomrule
\end{tabular}
\end{center}

\begin{itemize}
\tightlist
\item
  Competitive theory says employment should \textbf{fall} when the wage
  floor rises. The estimate is \(+2.76\) FTE per store, \(t \approx 2\).
\item
  The reaction split the profession. Whatever one believes about the
  economics, the paper put a \(2 \times 2\) in front of everyone, and
  the method's use in economics took off from here (Currie, Kleven and
  Zwiers 2020).
\item
  Note how a reader can check every number in this table. That is a
  feature.
\end{itemize}
\end{frame}

\begin{frame}{What the regression imputes}
\protect\phantomsection\label{what-the-regression-imputes}
\begin{center}
\includegraphics[width=1\linewidth,height=\textheight,keepaspectratio]{07a_did_fundamentals_files/figure-beamer/ck-diagram-1.pdf}
\end{center}

\begin{itemize}
\tightlist
\item
  \textbf{Left.} The regression draws NJ's missing November point by
  giving it PA's slope. The gap between NJ's actual November and that
  imputed point is \(\hat\delta = 2.76\).
\item
  \textbf{Right.} Same four observed points, same \(\hat\delta = 2.76\).
  But here NJ's true counterfactual fell by only \(0.5\), so the true
  ATT is \(1.09\). \textbf{The estimate does not know which picture it
  is in.}
\end{itemize}
\end{frame}

\section{5. Event Studies and
Falsifications}\label{event-studies-and-falsifications}

\begin{frame}{The event-study regression}
\protect\phantomsection\label{the-event-study-regression}
The \(2 \times 2\) interacts treatment with \emph{one} post dummy.
Interact it with \textbf{every period} except the baseline instead: \[
Y_{it} \;=\; \alpha_i + \lambda_t + \sum_{\tau \neq -1} \delta_\tau\, \big(D_i \times \mathbf{1}\{t - t^{*} = \tau\}\big) + \varepsilon_{it}
\]

\begin{itemize}
\tightlist
\item
  \(t^{*}\) is the treatment date, \(\tau\) the distance from it. The
  baseline \(\tau = -1\) is omitted: it must be untreated, which is no
  anticipation again.
\item
  \textbf{Each \(\hat\delta_\tau\) is its own \(2 \times 2\)} against
  the baseline. For \(\tau = -3\): \[
  \hat\delta_{-3} = \big(\bar Y_k^{\,\tau=-3} - \bar Y_k^{\,\tau=-1}\big) - \big(\bar Y_U^{\,\tau=-3} - \bar Y_U^{\,\tau=-1}\big)
  \]
\item
  Sometimes called ``long differences'': every period is compared to the
  same fixed baseline, however far away.
\end{itemize}
\end{frame}

\begin{frame}{What the coefficients measure}
\protect\phantomsection\label{what-the-coefficients-measure}
Substitute potential outcomes. Before treatment, no anticipation makes
\(\text{ATT}_{\tau} = 0\), so \[
\hat\delta_{-3} \;=\; \big(\mathbb{E}[Y^0_k|\tau=-3] - \mathbb{E}[Y^0_k|\tau=-1]\big) - \big(\mathbb{E}[Y^0_U|\tau=-3] - \mathbb{E}[Y^0_U|\tau=-1]\big)
\]

\begin{itemize}
\tightlist
\item
  \textbf{All four terms are observable.} A pre-period coefficient is
  the \emph{differential pre-trend} in untreated outcomes --- which is a
  fact about the data, \textbf{not the parallel trends assumption.}
  Parallel trends is about the post period, where one term is
  counterfactual.
\item
  \(\hat\delta_{-3} = +1.5\) means: from \(\tau=-3\) to the baseline,
  the treated group's outcome rose \(1.5\) more than the control's. Say
  it that way.
\item
  \textbf{After treatment}, \(\hat\delta_{+2}\) is the ATT at
  \(\tau = +2\) --- \emph{if} parallel trends holds from the baseline to
  \(\tau=+2\), there is no anticipation, and the control is never- or
  not-yet-treated. Tedious to say; if it feels tedious, you are not yet
  sure what you estimated.
\end{itemize}
\end{frame}

\begin{frame}{Castle doctrine: the raw data first}
\protect\phantomsection\label{castle-doctrine-the-raw-data-first}
\begin{center}
\includegraphics[width=0.8\linewidth,height=\textheight,keepaspectratio]{07a_did_fundamentals_files/figure-beamer/castle-raw-1.pdf}
\end{center}

\begin{itemize}
\tightlist
\item
  Cheng and Hoekstra plotted exactly this, one panel per cohort, so
  readers could compare the two paths with their own eyes. There is a
  lot to be said for that transparency.
\item
  Two flat, roughly parallel lines for six years; then the adopters move
  up while the others drift down. Note also what the picture does
  \emph{not} show: how either line would have moved in the other state
  of the world.
\end{itemize}
\end{frame}

\begin{frame}{Castle doctrine: the coefficients}
\protect\phantomsection\label{castle-doctrine-the-coefficients}
\begin{center}
\includegraphics[width=0.8\linewidth,height=\textheight,keepaspectratio]{07a_did_fundamentals_files/figure-beamer/castle-es-1.pdf}
\end{center}

\begin{itemize}
\tightlist
\item
  Each point is a \(2 \times 2\) against \(2005\), with a \(95\%\)
  interval, standard errors clustered by state. The hollow point is the
  omitted baseline: nothing was estimated there.
\item
  Pre-period coefficients: \(0.02\) to \(0.06\), none distinguishable
  from zero. Post: \(0.11\) in \(2006\), \(0.16\) in \(2007\), and
  positive throughout.
\item
  The \(2 \times 2\) of \(0.068\) is a weighted average of the five post
  coefficients, less a weighted average of the pre ones.
\end{itemize}
\end{frame}

\begin{frame}{Presenting an event study}
\protect\phantomsection\label{presenting-an-event-study}
\begin{itemize}
\tightlist
\item
  Plot the coefficients with their intervals. Two conventions:
  \textbf{disconnected} points with whiskers, or a \textbf{connected}
  line with a shaded band.
\item
  A case for disconnected: the band seems to pinch to zero at the
  baseline, which suggests precision where nothing was estimated; and
  you can count the coefficients.
\item
  Neither is correct. Choose what your audience reads correctly --- and
  make it beautiful enough that they want to look.
\end{itemize}

\vspace{0.4em}

\begin{itemize}
\tightlist
\item
  What a flat pre-period buys you is \textbf{suggestive} evidence, not a
  test. Parallel \emph{pre}-trends are not required for identification,
  and parallel \emph{post}-trends cannot be checked. Section 6 shows a
  case where the pre-trends break and parallel trends holds, and one
  where the pre-trends look fine and it does not.
\end{itemize}
\end{frame}

\begin{frame}{Falsification: the logic}
\protect\phantomsection\label{falsification-the-logic}
\begin{itemize}
\tightlist
\item
  A theory earns credibility by \textbf{excluding alternatives} (Popper
  1959). Your DiD result has two explanations: the treatment did it, or
  something else did.
\item
  A falsification takes the \emph{something else} seriously: if it were
  true, what else would we see? Then go and look.
\item
  Two designs:

  \begin{itemize}
  \tightlist
  \item
    \textbf{Same outcome, alternative group} --- a group the treatment
    cannot touch, but the alternative story would.
  \item
    \textbf{Same group, alternative outcome} --- an outcome the
    treatment cannot touch, but the alternative story would.
  \end{itemize}
\end{itemize}

\vspace{0.3em}

\begin{itemize}
\tightlist
\item
  Passing a falsification does not prove parallel trends; failing one
  does not prove the main result wrong. Both are \textbf{evidence}, and
  they weigh more than another asterisk.
\end{itemize}
\end{frame}

\begin{frame}{Falsification: two examples}
\protect\phantomsection\label{falsification-two-examples}
\begin{itemize}
\tightlist
\item
  \textbf{Same outcome, alternative group.} You find the minimum wage
  reduced low-wage employment. Alternative: the whole state economy was
  sinking. Then \textbf{high-wage} workers --- untouched by the floor
  --- should have lost jobs too. Re-run the \(2 \times 2\) on them. If
  they did, your main result has a rival explanation.
\item
  \textbf{Same group, alternative outcome.} Cigarette taxes reduce
  heart-attack admissions. Alternative: public insurance expanded at the
  same time, in the same states. Then admissions for
  \textbf{appendicitis and fractures} --- nothing to do with smoking ---
  should have fallen as well. If they did not, the rival loses.
\end{itemize}

\vspace{0.3em}

\begin{itemize}
\tightlist
\item
  The design needs a rival theory that makes a \textbf{prediction the
  treatment does not}. Without one, a ``placebo'' is just another
  regression.
\end{itemize}
\end{frame}

\begin{frame}{Falsification: milk is addictive and acne is contagious}
\protect\phantomsection\label{falsification-milk-is-addictive-and-acne-is-contagious}
\begin{itemize}
\tightlist
\item
  \textbf{Rational addiction} (Becker and Murphy 1988) gave economists
  an empirical model, which found alcohol, tobacco and gambling to be
  rationally addictive. Auld and Grootendorst (2004) ran the same model
  on \textbf{milk, eggs and oranges}. Milk came out among the most
  addictive goods in the study.
\item
  \textbf{Peer effects.} Obesity, smoking and happiness all seemed to
  spread through friendship networks (Christakis and Fowler 2007).
  Cohen-Cole and Fletcher (2008) ran the same model on \textbf{acne,
  height and headaches}. All three were contagious too.
\end{itemize}

\vspace{0.3em}

\begin{itemize}
\tightlist
\item
  Neither study attacked the theory. Each showed the \emph{empirical
  model} produces its result on data where the result cannot be true ---
  so its result on the real data is not evidence.
\item
  Tall people have tall friends because they are on the basketball team
  together. Ask what the equivalent is for your firms.
\end{itemize}
\end{frame}

\begin{frame}{Inference}
\protect\phantomsection\label{inference}
\begin{itemize}
\tightlist
\item
  Most DiD panels have \textbf{many periods}, treatment that varies
  \textbf{at the group level} (a state, an industry), and outcomes that
  are \textbf{serially correlated}. Bertrand, Duflo and Mullainathan
  (2004) showed that conventional standard errors then over-reject
  badly: placebo laws came out ``significant'' nearly half the time at
  the \(5\%\) level.
\item
  Three fixes, from their paper:

  \begin{enumerate}
  \tightlist
  \item
    \textbf{Cluster} at the level treatment is assigned --- the state,
    not the restaurant. Standard, and what every table in this lecture
    does.
  \item
    \textbf{Aggregate} to one pre and one post observation per unit: the
    long difference. Serial correlation has nowhere to live.
  \item
    \textbf{Block bootstrap}: resample whole states.
  \end{enumerate}
\end{itemize}

\vspace{0.3em}

\begin{itemize}
\tightlist
\item
  With \textbf{few clusters}, or few \emph{treated} clusters, clustering
  itself over-rejects. Castle doctrine has \(13\) adopters. Lecture 6a's
  wild cluster bootstrap applies; with a single treated unit,
  randomization inference (Buchmueller, DiNardo and Valletta 2011) is
  the honest option.
\end{itemize}
\end{frame}

\begin{frame}{DiD in the courtroom}
\protect\phantomsection\label{did-in-the-courtroom}
A prosecutor has three things: a \textbf{claim} of guilt,
\textbf{evidence} for it, and a \textbf{motive}. The claim is not
evidence, and a case with no motive is hard to believe beyond reasonable
doubt.

\begin{itemize}
\tightlist
\item
  Your regression table is the \textbf{claim}. A table of coefficients
  with asterisks is not a case.
\item
  Five elements of a strong DiD paper:

  \begin{enumerate}
  \tightlist
  \item
    \textbf{Bite}: the policy moved what it was meant to move.
  \item
    \textbf{Event studies}: the two groups moved together before.
  \item
    \textbf{Falsifications}: the rival explanations were given a chance
    to show up, and did not.
  \item
    \textbf{Main results}.
  \item
    \textbf{Mechanism}: \emph{why} the effect appears --- the motive.
  \end{enumerate}
\end{itemize}
\end{frame}

\begin{frame}{Medicaid and mortality (Miller, Johnson and Wherry 2021)}
\protect\phantomsection\label{medicaid-and-mortality-miller-johnson-and-wherry-2021}
The 2014 Medicaid expansion was adopted by some states and not others.
Outcome: annual mortality of \textbf{near-elderly} adults (55--64), from
death records linked to a large federal survey.

\begin{itemize}
\tightlist
\item
  \textbf{Bite} (three event studies): eligibility up, enrollment up,
  the uninsured share \emph{down} --- so enrollment was not just people
  switching from private cover.
\item
  \textbf{Falsification, same outcome, alternative group:} adults
  \textbf{65 and over} are on Medicare and cannot be affected. Their
  enrollment and mortality: no change.
\item
  \textbf{Main result:} mortality falls \(0.13\) percentage points,
  \(9.3\%\) of the mean, growing over time. The intervals are wide; the
  post coefficients are each different from zero but not from each
  other, and a linear trend cannot be rejected.
\item
  \textbf{Mechanism:} the decline is in disease-related deaths ---
  conditions that treatment, once you have access to it, prevents.
\end{itemize}

\vspace{0.2em}

\begin{itemize}
\tightlist
\item
  Could still be wrong: the counterfactual is unobserved. But the
  layered evidence leaves a sceptic with few places to stand. Read it
  for how the argument is built.
\end{itemize}
\end{frame}

\section{6. Who Gets Treated, and Why}\label{who-gets-treated-and-why}

\begin{frame}{The assignment mechanism}
\protect\phantomsection\label{the-assignment-mechanism}
\begin{itemize}
\tightlist
\item
  So far: an equation and an assumption. Now the question that decides
  whether the assumption is plausible: \textbf{why} were these units
  treated and not those?
\item
  Marx, Tamer and Tang (2024) and Ghanem, Sant'Anna and Wüthrich (2024)
  ask which assignment mechanisms are compatible with parallel trends.
  The answers are not what intuition suggests:
\end{itemize}

\begin{center}\footnotesize
\begin{tabular}{R{6.8cm}R{5.4cm}}
\toprule
Does not, by itself, break parallel trends & Breaks it \\
\midrule
a common trend in $Y^0$ & selection on the \textbf{gain} $Y^1 - Y^0$: \\
selection on a \textbf{fixed effect} & \quad the Roy model, the ``Perfect Doctor'' \\
selection on the \textbf{baseline level} of $Y^0$ (with a caveat) & selection on a \textbf{transitory dip} in $Y^0$ \\
selection on observables (conditional PT, next lecture) & \\
imperfect foresight about the future & \\
\bottomrule
\end{tabular}
\end{center}

\begin{itemize}
\tightlist
\item
  This knowledge comes from \textbf{outside the data}. Kahn-Lang and
  Lang (2019): say why the groups differed at baseline, and why that
  does not touch their trends.
\end{itemize}
\end{frame}

\begin{frame}{Good news: selection on fixed effects}
\protect\phantomsection\label{good-news-selection-on-fixed-effects}
\begin{center}
\includegraphics[width=0.74\linewidth,height=\textheight,keepaspectratio]{07a_did_fundamentals_files/figure-beamer/sel-fe-plot-1.pdf}
\end{center}

\(2{,}000\) workers, \(8\) periods, treated from period \(5\) with an
effect that grows by \(5\) each period (\(\times\) marks the truth).
\(Y^0\) has a worker effect, a common trend and an idiosyncratic part.
\textbf{Treatment goes to the workers with the highest fixed effects.}

\begin{itemize}
\tightlist
\item
  Permanent traits --- \emph{``there will always be a part of me that's
  sloppy and dirty''} --- are differenced out. Pre-coefficients at zero,
  post-coefficients on the truth.
\item
  Selection on \emph{who you are} is harmless. Selection on \emph{what
  is about to happen to you} is not.
\end{itemize}
\end{frame}

\begin{frame}{Good news with a caveat: selection on the baseline level}
\protect\phantomsection\label{good-news-with-a-caveat-selection-on-the-baseline-level}
\begin{center}
\includegraphics[width=0.7\linewidth,height=\textheight,keepaspectratio]{07a_did_fundamentals_files/figure-beamer/sel-y0-plot-1.pdf}
\end{center}

Treatment now goes to whoever is \textbf{lowest at the baseline}: a
means test, a programme for the worst schools.

\begin{itemize}
\tightlist
\item
  \textbf{Post-coefficients still on the truth.} Ghanem, Sant'Anna and
  Wüthrich (2024): selecting on the baseline level does not by itself
  break parallel trends.
\item
  \textbf{Pre-coefficients are not zero.} The baseline is both the
  selection point and the omitted period: the treated were picked for
  being low \emph{then}, so every other period looks high by comparison.
  A pre-trend test rejects, wrongly.
\item
  \textbf{The caveat.} True only if a low baseline says nothing about
  the \emph{future} change in \(Y^0\) (here, a random walk). If it is a
  \textbf{transitory dip} that reverts on its own, parallel trends fails
  too --- Ashenfelter's dip. Ask the institution, not the data.
\end{itemize}
\end{frame}

\begin{frame}{Bad news: the Perfect Doctor}
\protect\phantomsection\label{bad-news-the-perfect-doctor}
\begin{center}
\includegraphics[width=0.96\linewidth,height=\textheight,keepaspectratio]{07a_did_fundamentals_files/figure-beamer/pd-plot-1.pdf}
\end{center}

Gains \(\tau_i\) depend on age and GPA; \textbf{units are treated
exactly when \(\tau_i > 0\)} (Roy 1951; Heckman, Urzua and Vytlacil
2006). The \emph{return} to age and GPA in \(Y^0\) rises over time.

\begin{itemize}
\tightlist
\item
  Selecting on the gain selects on age and GPA, and those traits sit on
  a steeper \(Y^0\) path. \textbf{The groups were never going to move in
  parallel.} Here the pre-trends warn you; they need not.
\item
  This is the finance case: firms adopt, enter and list when they expect
  to gain. Unless assignment is at a level firms do not control --- a
  state law, a regulator's threshold, an index rule --- the Perfect
  Doctor is the default story.
\end{itemize}
\end{frame}

\begin{frame}{Difference-in-differences in finance}
\protect\phantomsection\label{difference-in-differences-in-finance}
\begin{itemize}
\tightlist
\item
  \textbf{Where the \(2 \times 2\) comes from.} State laws (Bertrand and
  Mullainathan 2003: anti-takeover laws and the ``quiet life'');
  regulatory thresholds and exemptions; exchange and index rules; court
  rulings; natural disasters. Treatment is assigned to firms by
  \emph{geography or rule}, not by choice.
\item
  \textbf{Bite.} Did the law actually bind? Which firms were exposed? If
  exposure is a choice, you are back to the Perfect Doctor.
\item
  \textbf{Clusters.} Treatment at the state level means clustering at
  the state level, and often few treated states. Bertrand, Duflo and
  Mullainathan was written about exactly these designs.
\item
  \textbf{Falsification.} Firms incorporated elsewhere but headquartered
  in the state; outcomes the law cannot reach; a placebo year.
\end{itemize}

\vspace{0.3em}

\begin{itemize}
\tightlist
\item
  And the question from this lecture: \textbf{why these states, these
  years?} A law passed because a state's firms were already struggling
  is selection on the trend.
\end{itemize}
\end{frame}

\begin{frame}{Conclusion}
\protect\phantomsection\label{conclusion}
\begin{enumerate}
\tightlist
\item
  \textbf{DiD is four averages and three subtractions.} Three
  regressions compute it exactly; regression adds standard errors, not
  identification.
\item
  \textbf{It identifies the ATT in the post period, plus a bias}: the
  \(2 \times 2\) on \(Y^0\). Parallel trends is the assumption that this
  bias is zero. It is about a counterfactual and \textbf{cannot be
  tested.}
\item
  \textbf{No anticipation} means the baseline is untreated. Violate it
  and you estimate the \emph{change} in the effect.
  \textbf{Already-treated controls} subtract the change in \emph{their}
  effect.
\item
  \textbf{Event studies} show whether the groups moved together
  \emph{before}; falsifications give rival explanations a chance to
  appear. Both are evidence, neither is proof.
\item
  \textbf{The assignment mechanism decides.} Selection on fixed traits
  or baseline levels can be fine; selection on gains, or on a dip that
  will revert, is not. Learn \emph{why} units were treated before you
  estimate anything.
\end{enumerate}

\vspace{0.3em}

\begin{itemize}
\tightlist
\item
  \textbf{Next:} more than one treatment date, covariates, and what
  two-way fixed effects does when the \(2 \times 2\)'s silently
  multiply.
\end{itemize}
\end{frame}

\begin{frame}{References}
\protect\phantomsection\label{references}
\tiny

\begin{itemize}
\tightlist
\item
  Auld, M. C. and P. Grootendorst (2004). An Empirical Analysis of Milk
  Addiction. \emph{Journal of Health Economics} 23(6), 1117--1133.
\item
  Baker, A., B. Callaway, S. Cunningham, A. Goodman-Bacon and P.
  Sant'Anna (2025). Difference-in-Differences Designs: A Practitioner's
  Guide. Working paper.
\item
  Bertrand, M., E. Duflo and S. Mullainathan (2004). How Much Should We
  Trust Differences-in-Differences Estimates? \emph{QJE} 119(1),
  249--275.
\item
  Bertrand, M. and S. Mullainathan (2003). Enjoying the Quiet Life?
  Corporate Governance and Managerial Preferences. \emph{JPE} 111(5),
  1043--1075.
\item
  Buchmueller, T., J. DiNardo and R. Valletta (2011). The Effect of an
  Employer Health Insurance Mandate on Health Insurance Coverage and the
  Demand for Labor. \emph{AEJ: Economic Policy} 3(4), 25--51.
\item
  Card, D. and A. Krueger (1994). Minimum Wages and Employment: A Case
  Study of the Fast-Food Industry in New Jersey and Pennsylvania.
  \emph{AER} 84(4), 772--793.
\item
  Cheng, C. and M. Hoekstra (2013). Does Strengthening Self-Defense Law
  Deter Crime or Escalate Violence? \emph{Journal of Human Resources}
  48(3), 821--854.
\item
  Cohen-Cole, E. and J. Fletcher (2008). Detecting Implausible Social
  Network Effects in Acne, Height, and Headaches. \emph{BMJ} 337, a2533.
\item
  Cunningham, S. (2025). \emph{Causal Inference: The Remix}, ch.~9 ---
  the source for this lecture, including the castle-doctrine data and
  the simulations.
\item
  Currie, J., H. Kleven and E. Zwiers (2020). Technology and Big Data
  Are Changing Economics. \emph{AEA Papers and Proceedings} 110, 42--48.
\item
  Ghanem, D., P. Sant'Anna and K. Wüthrich (2024). Selection and
  Parallel Trends. Working paper.
\item
  Goodman-Bacon, A. (2021). Difference-in-Differences with Variation in
  Treatment Timing. \emph{Journal of Econometrics} 225(2), 254--277.
\item
  Heckman, J., S. Urzua and E. Vytlacil (2006). Understanding
  Instrumental Variables in Models with Essential Heterogeneity.
  \emph{REStat} 88(3), 389--432.
\item
  Kahn-Lang, A. and K. Lang (2019). The Promise and Pitfalls of
  Differences-in-Differences. \emph{JBES} 38(3), 613--620.
\item
  Marx, P., E. Tamer and X. Tang (2024). Parallel Trends and Dynamic
  Choices. \emph{Journal of Political Economy Microeconomics} 2(1),
  129--171.
\item
  Miller, S., N. Johnson and L. Wherry (2021). Medicaid and Mortality:
  New Evidence from Linked Survey and Administrative Data. \emph{QJE}
  136(3), 1783--1829.
\item
  Rambachan, A. and J. Roth (2023). A More Credible Approach to Parallel
  Trends. \emph{Review of Economic Studies} 90(5), 2555--2591.
\item
  Roy, A. D. (1951). Some Thoughts on the Distribution of Earnings.
  \emph{Oxford Economic Papers} 3(2), 135--146.
\item
  Semmelweis, I. (1861). \emph{The Etiology, Concept, and Prophylaxis of
  Childbed Fever.} Clinic data: Wikipedia, ``Historical mortality rates
  of puerperal fever.''
\item
  Snow, J. (1855). \emph{On the Mode of Communication of Cholera.}
\end{itemize}
\end{frame}

\section{Backup}\label{backup}

\begin{frame}[fragile]{Stata: the four equivalences and the
no-anticipation simulation}
\protect\phantomsection\label{stata-the-four-equivalences-and-the-no-anticipation-simulation}
\begin{Shaded}
\begin{Highlighting}[]
\KeywordTok{use} \KeywordTok{data}\NormalTok{/castle.dta, }\KeywordTok{clear}
\KeywordTok{drop} \KeywordTok{if} \FunctionTok{inlist}\NormalTok{(effyear, 2005, 2007, 2008, 2009)          }\CommentTok{// 2006 cohort + never treated}
\KeywordTok{gen}\NormalTok{ treat = effyear == 2006}
\KeywordTok{drop} \KeywordTok{post}
\KeywordTok{gen} \KeywordTok{post}\NormalTok{  = }\FunctionTok{year}\NormalTok{ \textgreater{}= 2006}
\KeywordTok{reg}\NormalTok{ l\_homicide treat\#\#post, }\KeywordTok{cluster}\NormalTok{(sid)                  }\CommentTok{// interaction OLS  = 0.0682359}
\NormalTok{xtset sid }\FunctionTok{year}
\KeywordTok{xtreg}\NormalTok{ l\_homicide c.treat\#c.}\KeywordTok{post}\NormalTok{ i.}\FunctionTok{year}\NormalTok{, }\KeywordTok{fe} \KeywordTok{cluster}\NormalTok{(sid)   }\CommentTok{// two{-}way FE       = 0.0682359}
\KeywordTok{collapse}\NormalTok{ (}\KeywordTok{mean}\NormalTok{) l\_homicide, }\KeywordTok{by}\NormalTok{(sid treat }\KeywordTok{post}\NormalTok{)}
\KeywordTok{reshape} \KeywordTok{wide}\NormalTok{ l\_homicide, i(sid treat) j(}\KeywordTok{post}\NormalTok{)}
\KeywordTok{gen}\NormalTok{ dbar = l\_homicide1 {-} l\_homicide0}
\KeywordTok{reg}\NormalTok{ dbar treat, }\KeywordTok{cluster}\NormalTok{(sid)                              }\CommentTok{// long difference  = 0.0682359}
\end{Highlighting}
\end{Shaded}

\begin{itemize}
\tightlist
\item
  Full file, including the no-anticipation simulation (constant and
  dynamic effects, clean and treated baselines):
  \texttt{07a\_did\_sim.do}.
\end{itemize}
\end{frame}




\end{document}
